% Preâmbulo compartilhado dos cadernos de exercícios UNIFACOL.
% Definir \CadernoDisciplina e, opcionalmente, \CadernoCorHex antes deste input.
\providecommand{\CadernoDisciplina}{Disciplina}
\providecommand{\CadernoCorHex}{17365D}
\usepackage[a4paper,margin=2.05cm,headheight=15pt]{geometry}
\usepackage[T1]{fontenc}
\usepackage[utf8]{inputenc}
\usepackage[brazil]{babel}
\usepackage{lmodern,microtype,amsmath,amssymb}
\usepackage{enumitem,booktabs,tabularx,array,longtable,adjustbox,multirow}
\usepackage{xcolor,hyperref,graphicx,fancyhdr}
\usepackage[most]{tcolorbox}
\usepackage{tikz}
\usetikzlibrary{arrows.meta,positioning,fit,calc,patterns,shapes.geometric}
\definecolor{disciplinecolor}{HTML}{\CadernoCorHex}
\definecolor{stimulusgrey}{HTML}{EEF2F5}
\definecolor{answerorange}{HTML}{FFF0DE}
\definecolor{answerframe}{HTML}{B85C00}
\definecolor{supportgrey}{HTML}{F7F7F7}
\hypersetup{colorlinks=true,linkcolor=disciplinecolor,urlcolor=disciplinecolor}
\newcolumntype{Y}{>{\raggedright\arraybackslash}X}
\newcolumntype{C}{>{\centering\arraybackslash}X}
\setlength{\parindent}{0pt}
\setlength{\parskip}{.42em}
\setlist[enumerate]{leftmargin=*,itemsep=.28em,topsep=.35em}
\pagestyle{fancy}
\fancyhf{}
\lhead{UNIFACOL --- \CadernoDisciplina}
\rhead{Caderno ENADE --- 2026.2}
\cfoot{\thepage}
\newcommand{\ItemHeader}[4]{%
  \par\medskip\noindent
  \colorbox{disciplinecolor}{\parbox{\dimexpr\textwidth-2\fboxsep\relax}{\color{white}\bfseries Questão #1\hfill #2}}%
  \par\smallskip\noindent\textbf{Competência:} #3\par
  \noindent\textbf{Suporte:} #4\par\smallskip
}
\newcommand{\TextoBase}{\noindent\colorbox{stimulusgrey}{\parbox{\dimexpr\textwidth-2\fboxsep\relax}{\textbf{TEXTO-BASE}}}\par\smallskip}
\newcommand{\Comando}{\par\smallskip\noindent\textbf{COMANDO.} }
\newcommand{\FonteDidatica}{\par\smallskip\noindent{\footnotesize\textit{Fonte: situação e dados elaborados para fins didáticos.}}}
\newcommand{\FonteExterna}[1]{\par\smallskip\noindent{\footnotesize\textit{Fonte: #1}}}
\newcommand{\EspacoDiscursiva}{\par\noindent\rule{\textwidth}{.2pt}\vspace{1.15cm}\par\noindent\rule{\textwidth}{.2pt}}
\newtcolorbox{AnswerBox}{enhanced,breakable,colback=answerorange,colframe=answerframe,boxrule=.7pt,arc=1mm,title=Resposta explicada,fonttitle=\bfseries,coltitle=white,before skip=7pt,after skip=7pt}
\newtcolorbox{DocumentBox}[1][Documento ou artefato]{enhanced,breakable,colback=supportgrey,colframe=black!55,boxrule=.6pt,arc=.7mm,title={#1},fonttitle=\bfseries,coltitle=white,before skip=7pt,after skip=7pt}
\newtcolorbox{MemoBox}{enhanced,breakable,colback=white,colframe=disciplinecolor,boxrule=.7pt,arc=.7mm,title=Memorando,fonttitle=\bfseries,coltitle=white,before skip=7pt,after skip=7pt}
\newcommand{\LegendaDidatica}[1]{\par\smallskip{\footnotesize\textit{#1. Fonte: elaborado para fins didáticos.}}}
